\documentclass[11pt]{article}

\usepackage[utf8]{inputenc}
\usepackage[T1]{fontenc}
\usepackage[margin=1in]{geometry}
\usepackage{graphicx}
\usepackage{amsmath,amssymb}
\usepackage{booktabs}
\usepackage{array}
\usepackage{enumitem}
\usepackage{tcolorbox}
\usepackage{float}
\usepackage{placeins}
\usepackage[hidelinks]{hyperref}
\usepackage{textcomp}

\graphicspath{{plots/}{./}}

\newcommand{\Analog}{\textsc{Analog}}
\newcommand{\Alternating}{\textsc{Alternating}}
\newcommand{\FullSpiking}{\textsc{Full Spiking}}
\newcommand{\HybridLate}{\textsc{Hybrid Late}}
\newcommand{\HybridEarly}{\textsc{Hybrid Early}}
\newcommand{\HybridMid}{\textsc{Hybrid Mid}}

\newtcolorbox{keypoint}{%
  colback=black!3, colframe=black!60, boxrule=0.6pt, arc=1pt,
  left=8pt, right=8pt, top=6pt, bottom=6pt
}

\title{\Large Results Report:\\
Representation-Level Analysis of Rotation Invariance\\
in Spiking and Analog $p4$-Equivariant CNNs}
\author{Sahil Singh\\
CMI Summer Research Internship\\
Supervisor: Prof. K.V. Subrahmanyam}
\date{August 2026}

\begin{document}

\maketitle

\begin{abstract}
This report presents multi-seed (seeds 0, 1, 2) experimental results for four of the experiments
proposed in the research plan ``Spike Dynamics for Learning Invariant Representations'': layer-wise
emergence of invariance (Experiment 5.1), rotation sensitivity of individual neurons across layers
(Experiment 5.3), rotation sensitivity through time (Experiment 5.4), temporal evolution of
invariance via accuracy at truncated $T$ (Experiment 5.7), and direct equivariance error
(Experiment 5.8). Results are reported for four spike-placement configurations sharing an identical
$p4$-equivariant backbone: \Analog\ (no spiking layers), \Alternating\ (spikes at layers 2, 4, 6),
\FullSpiking\ (all 7 layers spike), and \HybridLate\ (only layer 6 spikes). All notation follows the
research plan exactly, restated in Section~\ref{sec:notation} for self-containedness. The central
empirical finding is that the layer-wise and temporal signatures of rotation sensitivity are shaped
substantially by the $p4$-equivariant architecture itself, independently of spike placement, while
spike placement modulates \emph{how} and \emph{when} that sensitivity manifests, in ways that do not
always follow the direction one would naively predict.
\end{abstract}

\tableofcontents

\section{How to Read This Report}

This report has three parts. Section~\ref{sec:notation} restates, without alteration, the notation
defined in the research plan ``Spike Dynamics for Learning Invariant Representations'' --- every
symbol used later is defined here first, in the same form KV sir's document uses it, so this report
can be read on its own or side by side with the research plan. Section~\ref{sec:setup} restates the
experimental setup. Sections~\ref{sec:finding1}--\ref{sec:finding4} present the actual multi-seed
results, organized by which of the nine proposed experiments (Section 5 of the research plan) they
answer. Figures could not be embedded in this pass; every figure has a labelled placeholder box
giving its exact filename on the GPU server (\texttt{gpuserver.cmi.ac.in}), so the corresponding PNG
can be inserted later without losing track of which plot goes where.

\section{Notation for Internal Representations}
\label{sec:notation}

This section restates Section 4 of the research plan exactly, together with the additional
quantities (direct equivariance error, truncated-$T$ accuracy) needed for Sections~\ref{sec:finding3}
and~\ref{sec:finding4}. Every symbol that appears in a later plot or table is defined here before it
is used.

\subsection{Layers, Inputs, and Rotations}

The network consists of layers
\[
x \to L_1 \to L_2 \to \cdots \to L_L, \qquad L = 7,
\]
where $x$ is a $28 \times 28$ input digit image (class 3, 4, 8, or 9), shown upright during training.
Given a rotation angle $\theta$, $R_\theta x$ denotes the same image rotated by $\theta$. Layer index
$\ell$ runs from 1 (nearest the input) to $L = 7$ (nearest the classifier); all layer-wise plots in
this report place $\ell$ on the horizontal axis running shallow $\to$ deep.

All four configurations examined here (Section~\ref{sec:setup}) share the identical $p4$-equivariant
backbone; they differ only in which layers use spiking (LIF) neurons versus analog (ReLU)
activations.

\subsection{Layer Representations}

For an input image $x$, $h_\ell(x)$ denotes the internal representation after layer $\ell$. In a
conventional (analog) neural network, $h_\ell(x)$ is simply the vector of activations of all neurons
in layer $\ell$.

In a spiking neural network each neuron evolves over time, so there are several possible notions of
representation. Throughout this work, following the research plan, we use the \emph{accumulated
spike counts} over the simulation horizon $T$. Suppose layer $\ell$ contains $n_\ell$ neurons and
neuron $i$ produces the spike train
\[
S_i(1), S_i(2), \ldots, S_i(T).
\]
Its accumulated spike count is
\begin{equation}
a_i(x) \;=\; \sum_{t=1}^{T} S_i(t). \label{eq:1}
\end{equation}
The layer representation is therefore
\[
h_\ell(x) \;=\; \big(a_1(x), a_2(x), \ldots, a_{n_\ell}(x)\big),
\]
the vector of spike counts of all neurons in layer $\ell$. For analog layers, $h_\ell(x)$ is the
ordinary activation vector; these are treated as a degenerate spiking layer with $T = 1$ so that the
same notation applies uniformly across configurations.

\paragraph{Why this matters.} The central question of the whole project is whether $h_\ell(x)$
becomes increasingly invariant to rotation ($h_\ell(x) \approx h_\ell(R_\theta x)$) as $\ell$
increases. Everything in Section~\ref{sec:finding1} is built directly on this comparison.

\subsection{Neuron Responses}

For an individual neuron $i$, $a_i(x)$ denotes its response to input $x$. Unlike a conventional
neuron, an SNN neuron produces an entire time series consisting of membrane potentials
$U_i(1), \ldots, U_i(T)$ and spike outputs $S_i(1), \ldots, S_i(T)$. Several choices are possible for
defining the neuron's scalar activation (accumulated spike count, average firing rate, final
membrane potential, average membrane potential); since classification throughout this work is based
on accumulated spike counts, we use $a_i(x) = \sum_{t=1}^{T} S_i(t)$ (Eq.~\ref{eq:1}) uniformly, as
specified in the research plan.

\subsection{Representation Similarity}
\label{sec:repsim}

To study where invariant representations emerge, we compare the layer representations of an image
and its rotated version, $h_\ell(x)$ versus $h_\ell(R_\theta x)$, using three similarity measures.
High similarity indicates that layer $\ell$ has learned a representation that is insensitive to
rotation.

\paragraph{Cosine similarity.} Treats $h_\ell(x)$ and $h_\ell(R_\theta x)$ as vectors and measures the
cosine of the angle between them (1 = identical direction, 0 = unrelated, $-1$ = opposite). This is
the strictest of the three measures: it requires exact neuron-by-neuron agreement in both direction
and relative magnitude, with no tolerance for the network re-using its neurons differently (e.g.\
permuting which unit encodes which feature) while preserving the same information content.

\paragraph{Centered Kernel Alignment (CKA).} Compares the Gram matrix of pairwise similarities
between a batch of images' representations under $h_\ell(x)$ against the Gram matrix under
$h_\ell(R_\theta x)$. Because CKA is invariant to orthogonal transformations and isotropic scaling of
the neuron basis, it tolerates the network computing the same information in a different internal
coordinate system, and is standard in the representation-similarity literature for this reason.
\emph{Implementation note:} for layer 1, the naive neuron-space Gram matrix has dimension
$\sim\!27{,}040 \times 27{,}040$ and is infeasible to materialize; we use the mathematically
equivalent batch-side Gram matrix formulation instead, verified to agree exactly with the naive form
on small test cases.

\paragraph{Singular Vector Canonical Correlation Analysis (SVCCA).} Projects each representation onto
its top principal directions, then finds the best linear alignment between the two reduced spaces and
reports the resulting correlation. SVCCA tolerates \emph{any} invertible linear transform relating
the two representations, making it the most permissive of the three measures. \emph{Caution:} SVCCA
is known to report spuriously high similarity when the batch size is too small relative to the neuron
count. A minimum batch-to-neuron ratio of $10\times$ is enforced throughout; even so, independent
(genuinely dissimilar) data was empirically found to produce a baseline SVCCA score of roughly
0.2--0.4, not 0. SVCCA values should therefore be read \emph{relatively} --- larger is more similar
--- never as an absolute 0-to-1 invariance scale.

\begin{table}[H]
\centering
\begin{tabular}{@{}p{2.6cm}p{3.6cm}p{4.2cm}p{2cm}@{}}
\toprule
\textbf{Metric} & \textbf{Tolerates} & \textbf{Strict about} & \textbf{Absolute scale?} \\
\midrule
Cosine similarity & Nothing --- exact match required & Neuron identity, magnitude, direction & Yes \\[4pt]
CKA & Rotation/rescaling of neuron basis & Pairwise similarity structure & Mostly yes \\[4pt]
SVCCA & Any invertible linear transform & Top principal (signal) directions & No --- compare relatively \\
\bottomrule
\end{tabular}
\caption{The three representation-similarity measures used in Experiment 5.1.}
\label{tab:metrics}
\end{table}

(Section~\ref{sec:finding2}) identifies where invariant representations are formed, at the
resolution of individual neurons rather than whole layers.

\subsection{Rotation Sensitivity of Individual Neurons}
\label{sec:Ri}

Following the research plan exactly, for neuron $i$ we define the rotation sensitivity
\begin{equation}
\mathcal{R}_i \;=\; \mathrm{Var}_\theta\big(a_i(R_\theta x)\big), \label{eq:2}
\end{equation}
where the variance is computed over the collection of 24 rotation angles $\theta$. A neuron with
$\mathcal{R}_i \approx 0$ responds similarly to all rotations and is therefore approximately
rotation-invariant; a large $\mathcal{R}_i$ indicates the neuron is strongly orientation-selective.
Studying the distribution of $\mathcal{R}_i$ across layers (Section~\ref{sec:finding2}) identifies
where invariant representations are formed, at the resolution of individual neurons rather than whole
layers.

\subsection{Temporal Representations and Rotation Sensitivity Through Time}

Instead of collapsing spike trains into accumulated spike counts, we may consider the instantaneous
representation
\[
h_\ell(t, x) \;=\; \big(S_1(t), S_2(t), \ldots, S_{n_\ell}(t)\big),
\]
recording the spike outputs of every neuron in layer $\ell$ at timestep $t$ specifically. Comparing
$h_\ell(t, x)$ with $h_\ell(t, R_\theta x)$ for successive $t$ investigates whether temporal
integration progressively builds rotation-invariant representations. Correspondingly,
$\mathcal{R}_i(t)$ denotes the rotation sensitivity of Eq.~\ref{eq:2} computed using only the spike
output at timestep $t$ (or, in the aggregate plots of Section~\ref{sec:54}, the median over neurons
of this quantity per layer, plotted against $t$).

\subsection{Direct Equivariance Error}

Experiment 5.8 of the research plan measures equivariance directly at the network output rather than
through an internal-representation proxy:
\[
\big\| f(Rx) - \rho(R) f(x) \big\|.
\]
For the classification setting used throughout this work, $f$ is the network's output
class-probability vector and the group action on the output, $\rho(R)$, is trivial (rotating the
input should not permute the class labels), so this reduces to the relative error
\begin{equation}
\text{Equivariance error}(\theta) \;=\; \frac{\big\| f(R_\theta x) - f(x) \big\|}{\big\| f(x) \big\|}. \label{eq:3}
\end{equation}
A perfectly rotation-invariant classifier has equivariance error 0 at every $\theta$. The
$p4$-equivariant architecture used here \emph{guarantees} error exactly 0 at the four group angles
$\theta \in \{0^\circ, 90^\circ, 180^\circ, 270^\circ\}$ by construction (Section~\ref{sec:arch}), and
makes no guarantee at any other angle, since those rotations are not elements of the $p4$ group the
convolution kernels are built from. Experiment 5.8 tests precisely whether error stays low everywhere
(full continuous invariance --- not claimed by this architecture) or dips to 0 only at the four group
angles while rising in between (the theoretically expected pattern).

\subsection{Truncated-$T$ Accuracy}

Experiment 5.7 of the research plan asks whether invariance emerges progressively as spikes
accumulate, by evaluating representations after $T = 1, 5, 10, 20, 50, 100$ timesteps. In this report
we use $T' \in \{1, 5, 10, 20, 50\}$: a single model is trained once using the full $T = 50$
timesteps, and at evaluation time its own recorded spike train is truncated to the first $T'$
timesteps to decode a classification decision, without retraining. This measures how quickly the
network's rotation-robust accuracy structure becomes available as a function of accumulated spiking
time.

\section{Experimental Setup}
\label{sec:setup}

\begin{keypoint}
Every result below is computed on the 4-class rotated-MNIST task (digits 3, 4, 8, 9), trained on
upright ($\theta = 0^\circ$) images only and evaluated at $\theta \in \{0^\circ, 15^\circ, \ldots,
345^\circ\}$ (24 angles), averaged or shown individually over three independent training seeds
$\{0, 1, 2\}$, for each of four spike-placement configurations described in Section~\ref{sec:setup}.
\end{keypoint}

\subsection{Architecture: $p4$-Equivariant Backbone}
\label{sec:arch}

All four configurations share an identical 7-layer $p4$-equivariant CNN backbone (six valid
$3 \times 3$ $p4$-convolutions and one valid $4 \times 4$ $p4$-convolution, 10 channels, one max-pool
after layer 2, a group-pooling step collapsing the 4 orientation channels via max, and a final linear
classifier), matching Cohen \& Welling's rotated-MNIST P4CNN specification. The group $p4$ consists
of translations composed with $90^\circ$ rotations; the convolution kernels are constructed so that
rotating the input by a multiple of $90^\circ$ produces an output that is provably (not
approximately) rotated/permuted correspondingly. This gives exact built-in invariance at
$\theta \in \{0^\circ, 90^\circ, 180^\circ, 270^\circ\}$ by mathematical construction and no
guarantee elsewhere --- which is precisely the reason the evaluation sweep uses all 24 angles rather
than only the four group angles: the interesting question is what happens \emph{between} the
guaranteed points.

\paragraph{Observation 1 (BatchNorm breaks equivariance).} \textit{Standard BatchNorm assigns each of
the 4 orientation channels independent running statistics, which breaks the architecture's
equivariance guarantee once trained upright-only (verified: an untrained model is equivariant to
$\sim 10^{-6}$; after training with standard BatchNorm, equivariance error rises to $0.1$--$0.2$). A
custom \texttt{P4BatchNorm2d}, sharing statistics and affine parameters across the 4 orientation
channels, restores equivariance to $\sim 10^{-6}$ even after training. All results in this report use
\texttt{P4BatchNorm2d} throughout.}

\subsection{Configurations}

All four configurations have identical architecture and parameter count (24,744 parameters); they
differ only in which layers use spiking LIF neurons (via \texttt{snntorch}, fast-sigmoid surrogate
gradient) versus analog ReLU activations.

\begin{table}[H]
\centering
\begin{tabular}{@{}lll@{}}
\toprule
\textbf{Configuration} & \textbf{Spiking layers} & \textbf{Description} \\
\midrule
\Analog & none & Zero-spiking control / ceiling reference \\
\Alternating & 2, 4, 6 & Spiking interleaved with analog throughout depth \\
\FullSpiking & all 1--7 & Maximal spiking \\
\HybridLate & 6 only & Spiking confined to near the output \\
\bottomrule
\end{tabular}
\caption{The four spike-placement configurations analyzed in this report. \HybridEarly\ (layers
1--2) and \HybridMid\ (layers 3--5), both implemented, have not yet been trained under the corrected
setup and are not covered here.}
\label{tab:configs}
\end{table}

\subsection{Training and Evaluation}

\begin{itemize}[leftmargin=1.5em]
\item \textbf{Input encoding.} Poisson rate coding over $T = 50$ timesteps (direct/constant-current
encoding produced a degenerate fixed-rate oscillation at low $T$ and was abandoned).
\item \textbf{Task.} 4-class rotated-MNIST, digits $\{3, 4, 8, 9\}$ --- chosen because none is
rotationally self-symmetric under $90^\circ/180^\circ$.
\item \textbf{Training.} Upright ($0^\circ$) images only, Adam, learning rate $10^{-3}$, batch size
64, 20 epochs, best checkpoint selected by $0^\circ$ test accuracy.
\item \textbf{Evaluation.} Full 24-angle sweep, $\theta \in \{0^\circ, 15^\circ, \ldots, 345^\circ\}$,
held-out test split.
\item \textbf{Seeds.} Three independent training runs per configuration, seeds $\{0, 1, 2\}$.
\item \textbf{Hardware.} \texttt{gpuserver.cmi.ac.in}, $2\times$ RTX 2080 Ti.
\end{itemize}

\subsection{Relation to the Original Internship Report}

The original report's compute-constrained single-seed runs ($T = 30$, 5 epochs, CPU only) reported
two dramatic mechanistic findings: \FullSpiking\ flat at 51--57\% accuracy across all angles, and
\HybridLate\ collapsing to exact chance (25\%) at non-group angles. Under the corrected setup above,
neither replicates: \FullSpiking\ reaches $\sim 99\%$ upright and a mean of $\sim 94\%$ across all
angles; \HybridLate's worst-case accuracy across all seeds and angles is 80.9\%. Three candidate
explanations (the \texttt{P4BatchNorm2d} fix, the longer schedule, the higher $T$) have been
identified but not yet isolated by controlled ablation (Section~\ref{sec:limitations}). The
representation-level experiments in Sections~\ref{sec:finding1}--\ref{sec:finding4} are a separate,
complementary investigation of the same underlying question posed by the research plan (Section 1):
even where accuracy differences between configurations have mostly disappeared, does the
\emph{internal representation} still reveal a meaningful role for spike dynamics?

\section{Baseline: Multi-Seed Accuracy Comparison}
\label{sec:baseline}

Before the representation-level results, we restate the accuracy picture with all three seeds and
error bars, since this is the backdrop against which Sections~\ref{sec:finding1}--\ref{sec:finding4}
should be read.

\begin{figure}[H]
\centering
\includegraphics[width=0.85\textwidth]{n4_comparison_final.png}
\caption{Config comparison, 4-class rotation-angle sweep. Mean $\pm$ std over $n = 3$ seeds for each
of \Analog, \Alternating, \FullSpiking, \HybridLate. Chance level (25\%) shown for reference.
\texttt{n4\_comparison\_final.png}}
\end{figure}

With all three seeds and error bars, the four configurations' rotation-accuracy curves sit within a
few percentage points of one another at every angle, with substantial error-bar overlap at the
hardest angles ($45^\circ, 135^\circ, 225^\circ, 315^\circ$). This confirms that the original
report's two dramatic accuracy-level findings (Section~\ref{sec:setup}) do not replicate under the
corrected setup, and that with proper multi-seed averaging the remaining accuracy differences
between configurations are small relative to seed-to-seed variation.

\begin{figure}[H]
\centering
\includegraphics[width=0.85\textwidth]{n4_rotation_curve.png}\\[6pt]
\includegraphics[width=0.85\textwidth]{n4_rotation_curve_training_curves.png}
\caption{\Analog-only rotation curve, mean $\pm$ std over 3 seeds ($\{0, 1, 2\}$), alongside per-seed
training-loss and upright-accuracy curves.
\texttt{n4\_rotation\_curve.png} and \texttt{n4\_rotation\_curve\_training\_curves.png}}
\end{figure}

\FloatBarrier

\section{Finding 1 --- Layer-wise Emergence of Invariance (Experiment 5.1)}
\label{sec:finding1}

\subsection{Setup}

Following Section 5.1 of the research plan exactly: for an image $x$ and its rotated version
$R_\theta x$, we compute $h_\ell(x)$ and $h_\ell(R_\theta x)$ at every hidden layer
$\ell = 1, \ldots, 7$, and measure similarity using cosine similarity, CKA, and SVCCA
(Section~\ref{sec:repsim}), plotted against layer depth for each of the 24 test angles.

\subsection{Result: A Common Depth-Wise Collapse in Every Configuration}

Every configuration --- including \Analog, the zero-spiking control --- shows the same qualitative
shape: CKA similarity to the upright representation stays at or above roughly 0.97--0.99 through
layers 1--5, and only then drops sharply in layers 6--7, to roughly 0.70--0.85 at the hardest angles
(typically the non-group angles nearest a group angle: $15^\circ, 30^\circ, 45^\circ, 60^\circ$,
etc.). The last two layers are where representational invariance to rotation is lost; the first five
preserve it almost perfectly, in every configuration tested.

\paragraph{Aggregate CKA summary (one line per angle, per configuration)}

\begin{figure}[H]
\centering
\includegraphics[width=0.8\textwidth]{analog_layerwise_cka.png}
\caption{\Analog\ --- CKA similarity to upright vs.\ layer, one line per test angle.
\texttt{analog\_layerwise\_cka.png}}
\end{figure}

\begin{figure}[H]
\centering
\includegraphics[width=0.8\textwidth]{alternating_layerwise_cka.png}
\caption{\Alternating\ --- CKA similarity to upright vs.\ layer, one line per test angle.
\texttt{alternating\_layerwise\_cka.png}}
\end{figure}

\begin{figure}[H]
\centering
\includegraphics[width=0.8\textwidth]{full_layerwise_cka.png}
\caption{\FullSpiking\ --- CKA similarity to upright vs.\ layer, one line per test angle.
\texttt{full\_layerwise\_cka.png}}
\end{figure}

\noindent\HybridLate\ --- CKA similarity to upright vs.\ layer, one line per test angle.
\texttt{hybrid\_late\_layerwise\_cka.png} was not generated as a standalone aggregate for this pass;
see the per-seed Section~5.1 plots below for \HybridLate\ instead.

\FloatBarrier

\subsection{Per-Seed, All-Three-Metrics Detail}

The figures below show, for each configuration and seed, all three similarity measures side by side
with all 24 angles overlaid, confirming the layer-6/7 collapse is not an artifact of one metric and
is stable across seeds.

\paragraph{Analog}

\begin{figure}[H]
\centering
\includegraphics[width=0.8\textwidth]{analog_seed0/analog_seed0_5.1_layerwise_similarity.png}
\caption{\Analog, seed 0 --- cosine / CKA / SVCCA vs.\ layer, all 24 angles.
\texttt{analog\_seed0/analog\_seed0\_5.1\_layerwise\_similarity.png}}
\end{figure}

\begin{figure}[H]
\centering
\includegraphics[width=0.8\textwidth]{analog_seed1/analog_seed1_5.1_layerwise_similarity.png}
\caption{\Analog, seed 1.
\texttt{analog\_seed1/analog\_seed1\_5.1\_layerwise\_similarity.png}}
\end{figure}

\begin{figure}[H]
\centering
\includegraphics[width=0.8\textwidth]{analog_seed2/analog_seed2_5.1_layerwise_similarity.png}
\caption{\Analog, seed 2.
\texttt{analog\_seed2/analog\_seed2\_5.1\_layerwise\_similarity.png}}
\end{figure}

\FloatBarrier
\paragraph{Alternating (spikes at layers 2, 4, 6)}

\begin{figure}[H]
\centering
\includegraphics[width=0.8\textwidth]{alternating_seed0/alternating_seed0_5.1_layerwise_similarity.png}
\caption{\Alternating, seed 0.
\texttt{alternating\_seed0/alternating\_seed0\_5.1\_layerwise\_similarity.png}}
\end{figure}

\begin{figure}[H]
\centering
\includegraphics[width=0.8\textwidth]{alternating_seed1/alternating_seed1_5.1_layerwise_similarity.png}
\caption{\Alternating, seed 1.
\texttt{alternating\_seed1/alternating\_seed1\_5.1\_layerwise\_similarity.png}}
\end{figure}

\begin{figure}[H]
\centering
\includegraphics[width=0.8\textwidth]{alternating_seed2/alternating_seed2_5.1_layerwise_similarity.png}
\caption{\Alternating, seed 2.
\texttt{alternating\_seed2/alternating\_seed2\_5.1\_layerwise\_similarity.png}}
\end{figure}

\FloatBarrier
\paragraph{Full Spiking (all 7 layers spike)}

\begin{figure}[H]
\centering
\includegraphics[width=0.8\textwidth]{full_seed0/full_seed0_5.1_layerwise_similarity.png}
\caption{\FullSpiking, seed 0.
\texttt{full\_seed0/full\_seed0\_5.1\_layerwise\_similarity.png}}
\end{figure}

\begin{figure}[H]
\centering
\includegraphics[width=0.8\textwidth]{full_seed1/full_seed1_5.1_layerwise_similarity.png}
\caption{\FullSpiking, seed 1.
\texttt{full\_seed1/full\_seed1\_5.1\_layerwise\_similarity.png}}
\end{figure}

\begin{figure}[H]
\centering
\includegraphics[width=0.8\textwidth]{full_seed2/full_seed2_5.1_layerwise_similarity.png}
\caption{\FullSpiking, seed 2.
\texttt{full\_seed2/full\_seed2\_5.1\_layerwise\_similarity.png}}
\end{figure}

\FloatBarrier
\paragraph{Hybrid Late (layer 6 only spikes)}

\begin{figure}[H]
\centering
\includegraphics[width=0.8\textwidth]{hybrid_late_seed0/hybrid_late_seed0_5.1_layerwise_similarity.png}
\caption{\HybridLate, seed 0.
\texttt{hybrid\_late\_seed0/hybrid\_late\_seed0\_5.1\_layerwise\_similarity.png}}
\end{figure}

\begin{figure}[H]
\centering
\includegraphics[width=0.8\textwidth]{hybrid_late_seed1/hybrid_late_seed1_5.1_layerwise_similarity.png}
\caption{\HybridLate, seed 1.
\texttt{hybrid\_late\_seed1/hybrid\_late\_seed1\_5.1\_layerwise\_similarity.png}}
\end{figure}

\begin{figure}[H]
\centering
\includegraphics[width=0.8\textwidth]{hybrid_late_seed2/hybrid_late_seed2_5.1_layerwise_similarity.png}
\caption{\HybridLate, seed 2.
\texttt{hybrid\_late\_seed2/hybrid\_late\_seed2\_5.1\_layerwise\_similarity.png}}
\end{figure}

\FloatBarrier

\subsection{Reading the Three Metrics Together}

The three panels within each figure disagree with each other informatively rather than redundantly.
Cosine similarity is much lower and noisier throughout (typically 0.3--0.6 even at shallow layers,
falling to 0.2--0.3 by layer 7), reflecting its strictness: it demands exact neuron-by-neuron
agreement and penalizes even a meaningless internal reshuffling of which channel encodes which
feature, even when the information content is unchanged. CKA and SVCCA both stay near 1.0 through
layer 5 and only then diverge slightly at layers 6--7, with SVCCA sitting a little above CKA at the
deepest layers in several seeds --- consistent with SVCCA's more permissive tolerance for any linear
transform (Section~\ref{sec:repsim}) crediting structure that CKA's stricter pairwise-similarity
criterion does not. The three metrics agreeing on \emph{where} the drop happens while disagreeing on
\emph{how severe} it looks is itself informative: it indicates a real, robust change in what the
layer computes, not an artifact of one particular similarity definition.

\subsection{Interpretation}

\begin{enumerate}[leftmargin=1.8em]
\item The representational collapse at layers 6--7 occurs in \Analog\ too, with zero spiking neurons
anywhere in the network. Whatever causes deep layers to become rotation-sensitive is a property of
the $p4$-equivariant architecture itself --- most plausibly the group-pooling and final
linear-classifier stages, the only components not built from strictly equivariant convolutions ---
not something introduced by spiking dynamics.
\item Spike placement modulates the depth pattern's shape somewhat but not its existence or
location. \Alternating\ and \FullSpiking\ show a slightly earlier onset of divergence (visible
already around layer 5) relative to \Analog\ and \HybridLate, but all four configurations land in
the same 0.70--0.85 CKA range by layer 7 at the hardest angles.
\item This complicates the research plan's central hypothesis (Section 6 of the research plan,
restated in Section~\ref{sec:synthesis} of this report) in an interesting way. The hypothesis space
as originally framed treats spike dynamics as the primary lever for building (or failing to build)
invariant representations. Finding 1 instead shows that the \emph{architecture} sets both the ceiling
and the location of where invariance is lost, and spike dynamics act as a secondary modulator on top
of that. This motivates reframing the guiding question from ``do spikes build invariance?'' to
``given that architecture alone determines where invariance breaks down, what does spike placement
change about how that breakdown proceeds?'' --- precisely what Findings 2 and 3 investigate.
\end{enumerate}

\section{Finding 2 --- Rotation Sensitivity of Individual Neurons, by Layer and by Time (Experiments 5.3, 5.4)}
\label{sec:finding2}

Finding 1 examined whole-layer representations. This section zooms into individual neurons using
$\mathcal{R}_i$ (Eq.~\ref{eq:2}, Section~\ref{sec:Ri}), asking two complementary questions: which
layers contain the most rotation-sensitive individual neurons (Experiment 5.3), and does
$\mathcal{R}_i(t)$ decrease as more spiking timesteps are accumulated, as the research plan's
temporal-averaging hypothesis (Section 9.4 of the research plan) would predict (Experiment 5.4)?

\subsection{5.3 --- Distribution of $\mathcal{R}_i$ Across Layers}

Each boxplot shows, for one configuration, the distribution of $\mathcal{R}_i$ across all neurons in
each of the 7 layers (log scale, since $\mathcal{R}_i$ spans several orders of magnitude). The orange
line is the median; the box spans the interquartile range; whiskers extend to the typical range with
outliers beyond.

\paragraph{Aggregate boxplots (seeds pooled per configuration)}

\begin{figure}[H]
\centering
\includegraphics[width=0.8\textwidth]{analog_rotation_sensitivity.png}
\caption{\Analog\ --- $\mathcal{R}_i$ distribution by layer.
\texttt{analog\_rotation\_sensitivity.png}}
\end{figure}

\begin{figure}[H]
\centering
\includegraphics[width=0.8\textwidth]{alternating_rotation_sensitivity.png}
\caption{\Alternating\ --- $\mathcal{R}_i$ distribution by layer.
\texttt{alternating\_rotation\_sensitivity.png}}
\end{figure}

\begin{figure}[H]
\centering
\includegraphics[width=0.8\textwidth]{full_rotation_sensitivity.png}
\caption{\FullSpiking\ --- $\mathcal{R}_i$ distribution by layer.
\texttt{full\_rotation\_sensitivity.png}}
\end{figure}

\noindent\HybridLate\ --- $\mathcal{R}_i$ distribution by layer.
\texttt{hybrid\_late\_rotation\_sensitivity.png} was not generated as a standalone aggregate for this
pass; see the per-seed plots below for \HybridLate\ instead.

\FloatBarrier

\paragraph{Per-seed boxplots}

\begin{figure}[H]
\centering
\includegraphics[width=0.8\textwidth]{analog_seed0/analog_seed0_5.3_rotation_sensitivity.png}
\caption{\Analog, seed 0.
\texttt{analog\_seed0/analog\_seed0\_5.3\_rotation\_sensitivity.png}}
\end{figure}

\begin{figure}[H]
\centering
\includegraphics[width=0.8\textwidth]{analog_seed1/analog_seed1_5.3_rotation_sensitivity.png}
\caption{\Analog, seed 1.
\texttt{analog\_seed1/analog\_seed1\_5.3\_rotation\_sensitivity.png}}
\end{figure}

\begin{figure}[H]
\centering
\includegraphics[width=0.8\textwidth]{analog_seed2/analog_seed2_5.3_rotation_sensitivity.png}
\caption{\Analog, seed 2.
\texttt{analog\_seed2/analog\_seed2\_5.3\_rotation\_sensitivity.png}}
\end{figure}

\begin{figure}[H]
\centering
\includegraphics[width=0.8\textwidth]{alternating_seed0/alternating_seed0_5.3_rotation_sensitivity.png}
\caption{\Alternating, seed 0.
\texttt{alternating\_seed0/alternating\_seed0\_5.3\_rotation\_sensitivity.png}}
\end{figure}

\begin{figure}[H]
\centering
\includegraphics[width=0.8\textwidth]{alternating_seed1/alternating_seed1_5.3_rotation_sensitivity.png}
\caption{\Alternating, seed 1.
\texttt{alternating\_seed1/alternating\_seed1\_5.3\_rotation\_sensitivity.png}}
\end{figure}

\begin{figure}[H]
\centering
\includegraphics[width=0.8\textwidth]{alternating_seed2/alternating_seed2_5.3_rotation_sensitivity.png}
\caption{\Alternating, seed 2.
\texttt{alternating\_seed2/alternating\_seed2\_5.3\_rotation\_sensitivity.png}}
\end{figure}

\begin{figure}[H]
\centering
\includegraphics[width=0.8\textwidth]{full_seed0/full_seed0_5.3_rotation_sensitivity.png}
\caption{\FullSpiking, seed 0.
\texttt{full\_seed0/full\_seed0\_5.3\_rotation\_sensitivity.png}}
\end{figure}

\begin{figure}[H]
\centering
\includegraphics[width=0.8\textwidth]{full_seed1/full_seed1_5.3_rotation_sensitivity.png}
\caption{\FullSpiking, seed 1.
\texttt{full\_seed1/full\_seed1\_5.3\_rotation\_sensitivity.png}}
\end{figure}

\begin{figure}[H]
\centering
\includegraphics[width=0.8\textwidth]{full_seed2/full_seed2_5.3_rotation_sensitivity.png}
\caption{\FullSpiking, seed 2.
\texttt{full\_seed2/full\_seed2\_5.3\_rotation\_sensitivity.png}}
\end{figure}

\begin{figure}[H]
\centering
\includegraphics[width=0.8\textwidth]{hybrid_late_seed0/hybrid_late_seed0_5.3_rotation_sensitivity.png}
\caption{\HybridLate, seed 0.
\texttt{hybrid\_late\_seed0/hybrid\_late\_seed0\_5.3\_rotation\_sensitivity.png}}
\end{figure}

\begin{figure}[H]
\centering
\includegraphics[width=0.8\textwidth]{hybrid_late_seed1/hybrid_late_seed1_5.3_rotation_sensitivity.png}
\caption{\HybridLate, seed 1.
\texttt{hybrid\_late\_seed1/hybrid\_late\_seed1\_5.3\_rotation\_sensitivity.png}}
\end{figure}

\begin{figure}[H]
\centering
\includegraphics[width=0.8\textwidth]{hybrid_late_seed2/hybrid_late_seed2_5.3_rotation_sensitivity.png}
\caption{\HybridLate, seed 2.
\texttt{hybrid\_late\_seed2/hybrid\_late\_seed2\_5.3\_rotation\_sensitivity.png}}
\end{figure}

\FloatBarrier

\paragraph{What differs between configurations}

All four configurations place layer 7 at or near the top of the $\mathcal{R}_i$ ranking, consistent
with Finding 1's identification of the last layers as where invariance breaks down. The shape of the
climb to layer 7 differs meaningfully:

\begin{itemize}[leftmargin=1.5em]
\item \textbf{Alternating:} a sharp, order-of-magnitude jump only at layer 7 --- median
$\mathcal{R}_i \sim 10$--100 for layers 1--6, jumping to $\sim 400$--800 at layer 7.
\item \textbf{Full Spiking:} a gradual, full-depth climb --- median $\mathcal{R}_i$ increases roughly
monotonically from $\sim 10$ (layer 1) to $\sim 90$--100 (layer 7), with no single sharp jump.
\item \textbf{Hybrid Late:} nearly flat across all seven layers --- median $\mathcal{R}_i$ sits in a
comparatively narrow band ($\sim 10$--80) at every depth, including layer 7.
\item \textbf{Analog:} a gradual full-depth climb similar in shape to \FullSpiking, but at a
dramatically smaller absolute scale ($\sim 0.01$--0.05 throughout, roughly $1000\times$ smaller than
the spiking configurations at comparable layers).
\end{itemize}

\paragraph{Observation 2 ($\mathcal{R}_i$ scale is not comparable across activation types).}
\textit{Because $\mathcal{R}_i$'s absolute scale depends on whether the underlying activation is a
continuous ReLU output or a spike count, the fair cross-configuration comparison is the \emph{shape}
of the $\mathcal{R}_i$-vs-layer curve (sharp late jump vs.\ gradual climb vs.\ flat), not the raw
magnitude. Within a single configuration, comparisons across layers or across seeds are on safe
footing.}

\subsection{5.4 --- Rotation Sensitivity Through Time}
\label{sec:54}

This experiment is the direct empirical test of the research plan's Section 9.4 hypothesis: does
temporal averaging progressively stabilize the representation, so that
$\mathrm{Var}(\bar h_T) = O(T^{-1})$ and $\mathcal{R}_i(t)$ decays toward 0 as $t$ grows? The plots
below show median $\mathcal{R}_i(t)$ per layer against timestep $t \in \{1, \ldots, 50\}$, log scale
on the vertical axis.

\begin{figure}[H]
\centering
\includegraphics[width=0.8\textwidth]{alternating_seed0/alternating_seed0_5.4_rotation_sensitivity_through_time.png}
\caption{\Alternating, seed 0 --- $\mathcal{R}_i(t)$ vs.\ timestep, all 7 layers.
\texttt{alternating\_seed0/alternating\_seed0\_5.4\_rotation\_sensitivity\_through\_time.png}}
\end{figure}

\begin{figure}[H]
\centering
\includegraphics[width=0.8\textwidth]{alternating_seed1/alternating_seed1_5.4_rotation_sensitivity_through_time.png}
\caption{\Alternating, seed 1.
\texttt{alternating\_seed1/alternating\_seed1\_5.4\_rotation\_sensitivity\_through\_time.png}}
\end{figure}

\begin{figure}[H]
\centering
\includegraphics[width=0.8\textwidth]{alternating_seed2/alternating_seed2_5.4_rotation_sensitivity_through_time.png}
\caption{\Alternating, seed 2.
\texttt{alternating\_seed2/alternating\_seed2\_5.4\_rotation\_sensitivity\_through\_time.png}}
\end{figure}

\begin{figure}[H]
\centering
\includegraphics[width=0.8\textwidth]{full_seed0/full_seed0_5.4_rotation_sensitivity_through_time.png}
\caption{\FullSpiking, seed 0.
\texttt{full\_seed0/full\_seed0\_5.4\_rotation\_sensitivity\_through\_time.png}}
\end{figure}

\begin{figure}[H]
\centering
\includegraphics[width=0.8\textwidth]{full_seed1/full_seed1_5.4_rotation_sensitivity_through_time.png}
\caption{\FullSpiking, seed 1.
\texttt{full\_seed1/full\_seed1\_5.4\_rotation\_sensitivity\_through\_time.png}}
\end{figure}

\begin{figure}[H]
\centering
\includegraphics[width=0.8\textwidth]{full_seed2/full_seed2_5.4_rotation_sensitivity_through_time.png}
\caption{\FullSpiking, seed 2.
\texttt{full\_seed2/full\_seed2\_5.4\_rotation\_sensitivity\_through\_time.png}}
\end{figure}

\begin{figure}[H]
\centering
\includegraphics[width=0.8\textwidth]{hybrid_late_seed0/hybrid_late_seed0_5.4_rotation_sensitivity_through_time.png}
\caption{\HybridLate, seed 0.
\texttt{hybrid\_late\_seed0/hybrid\_late\_seed0\_5.4\_rotation\_sensitivity\_through\_time.png}}
\end{figure}

\begin{figure}[H]
\centering
\includegraphics[width=0.8\textwidth]{hybrid_late_seed1/hybrid_late_seed1_5.4_rotation_sensitivity_through_time.png}
\caption{\HybridLate, seed 1.
\texttt{hybrid\_late\_seed1/hybrid\_late\_seed1\_5.4\_rotation\_sensitivity\_through\_time.png}}
\end{figure}

\begin{figure}[H]
\centering
\includegraphics[width=0.8\textwidth]{hybrid_late_seed2/hybrid_late_seed2_5.4_rotation_sensitivity_through_time.png}
\caption{\HybridLate, seed 2.
\texttt{hybrid\_late\_seed2/hybrid\_late\_seed2\_5.4\_rotation\_sensitivity\_through\_time.png}}
\end{figure}

\FloatBarrier

\paragraph{Observation 3 (Rise-then-plateau, not decay).} \textit{In every configuration and every
seed, every layer's $\mathcal{R}_i(t)$ curve rises quickly over the first 2--5 timesteps (as the
membrane potential accumulates its first spikes and begins producing non-degenerate output), then
plateaus for the remaining $\sim 45$ timesteps with only small step-to-step jitter and no visible
downward trend. Layer 7 plateaus visibly higher than the other layers in \Alternating\ and
\FullSpiking\ (consistent with Section 5.1 of this report); in \HybridLate\ the plateau levels across
layers are closer together, consistent with that configuration's flatter layer-wise sensitivity
profile.}

\Analog\ has no meaningful 5.4 plot: with $T = 1$ there is no timestep axis to show evolving over
time.

This directly contradicts the ``temporal averaging reduces sensitivity'' hypothesis of
research-plan Section 9.4, as measured on these trained checkpoints: each layer's rotation
sensitivity is set almost immediately, within the first handful of timesteps, and then stays fixed
for the remainder of the simulation. Temporal integration in this architecture does not act as a
noise-averaging invariance-builder over the timescales examined here --- the theoretical contraction
property $d_{t+1} \leq \beta d_t$ proposed in research-plan Section 9.5 is not observed empirically
for $\mathcal{R}_i(t)$ once the initial transient has passed.

\section{Finding 3 --- Temporal Evolution of Invariance via Truncated-$T$ Accuracy (Experiment 5.7)}
\label{sec:finding3}

Section~\ref{sec:finding2} examined representation-level sensitivity through time. This section asks
the same question at the level that ultimately matters for a user of the classifier: accuracy. Using
the truncated-$T$ evaluation of Section 2.8, we decode a classification decision from the same
trained model using only its first $T' \in \{1, 5, 10, 20, 50\}$ timesteps, and plot accuracy against
all 24 test angles for each $T'$. This produces, for each $T'$, the familiar W-shaped
rotation-accuracy curve (peaks at the group angles, troughs between them), and lets us observe how
quickly that shape ``switches on'' as more timesteps become available --- directly answering the
question posed in research-plan Section 5.7: does invariance emerge progressively as spikes
accumulate?

\subsection{Full Spiking: Requires Real Accumulation Time}

\begin{figure}[H]
\centering
\includegraphics[width=0.8\textwidth]{full_seed0/full_seed0_5.7_temporal_evolution.png}
\caption{\FullSpiking, seed 0 --- accuracy vs.\ angle at $T' = 1, 5, 10, 20, 50$.
\texttt{full\_seed0/full\_seed0\_5.7\_temporal\_evolution.png}}
\end{figure}

\begin{figure}[H]
\centering
\includegraphics[width=0.8\textwidth]{full_seed1/full_seed1_5.7_temporal_evolution.png}
\caption{\FullSpiking, seed 1.
\texttt{full\_seed1/full\_seed1\_5.7\_temporal\_evolution.png}}
\end{figure}

\begin{figure}[H]
\centering
\includegraphics[width=0.8\textwidth]{full_seed2/full_seed2_5.7_temporal_evolution.png}
\caption{\FullSpiking, seed 2.
\texttt{full\_seed2/full\_seed2\_5.7\_temporal\_evolution.png}}
\end{figure}

\FloatBarrier

In all three seeds, the $T' = 1$ curve is flat at approximately chance level ($\sim 25$--27\%) across
every angle: a single timestep carries insufficient accumulated spike information for any confident
decision, correct or otherwise. Real W-shaped structure appears only by $T' = 10$, and is nearly
indistinguishable from the full $T' = 50$ curve by $T' = 20$. This is mechanistically sensible: with
all 7 layers spiking, every layer including the first needs several timesteps to integrate enough
input spikes to pass a meaningful signal onward, and this startup delay compounds across all 7
layers.

\subsection{Alternating: Partial Signal Already at $T' = 1$}

\begin{figure}[H]
\centering
\includegraphics[width=0.8\textwidth]{alternating_seed0/alternating_seed0_5.7_temporal_evolution.png}
\caption{\Alternating, seed 0 --- accuracy vs.\ angle at $T' = 1, 5, 10, 20, 50$.
\texttt{alternating\_seed0/alternating\_seed0\_5.7\_temporal\_evolution.png}}
\end{figure}

\begin{figure}[H]
\centering
\includegraphics[width=0.8\textwidth]{alternating_seed1/alternating_seed1_5.7_temporal_evolution.png}
\caption{\Alternating, seed 1.
\texttt{alternating\_seed1/alternating\_seed1\_5.7\_temporal\_evolution.png}}
\end{figure}

\begin{figure}[H]
\centering
\includegraphics[width=0.8\textwidth]{alternating_seed2/alternating_seed2_5.7_temporal_evolution.png}
\caption{\Alternating, seed 2.
\texttt{alternating\_seed2/alternating\_seed2\_5.7\_temporal\_evolution.png}}
\end{figure}

\FloatBarrier

Unlike \FullSpiking, \Alternating's $T' = 1$ curve is not flat: it shows a weak but genuine W-shape,
oscillating roughly between 25\% and 50--70\% (seed- and angle-dependent), tracking the peaks and
troughs of the full-$T$ curve well below its absolute accuracy. This is architecturally sensible:
with layers 1, 3, 5, 7 analog, the network has several ``instant'' processing stages interleaved with
the spiking ones, so even one timestep of spiking input still passes through enough analog
processing to carry rotation-relevant signal to the output.

\subsection{Hybrid Late: Nearly Full Structure Already at $T' = 1$}

\begin{figure}[H]
\centering
\includegraphics[width=0.8\textwidth]{hybrid_late_seed0/hybrid_late_seed0_5.7_temporal_evolution.png}
\caption{\HybridLate, seed 0 --- accuracy vs.\ angle at $T' = 1, 5, 10, 20, 50$.
\texttt{hybrid\_late\_seed0/hybrid\_late\_seed0\_5.7\_temporal\_evolution.png}}
\end{figure}

\begin{figure}[H]
\centering
\includegraphics[width=0.8\textwidth]{hybrid_late_seed1/hybrid_late_seed1_5.7_temporal_evolution.png}
\caption{\HybridLate, seed 1.
\texttt{hybrid\_late\_seed1/hybrid\_late\_seed1\_5.7\_temporal\_evolution.png}}
\end{figure}

\begin{figure}[H]
\centering
\includegraphics[width=0.8\textwidth]{hybrid_late_seed2/hybrid_late_seed2_5.7_temporal_evolution.png}
\caption{\HybridLate, seed 2.
\texttt{hybrid\_late\_seed2/hybrid\_late\_seed2\_5.7\_temporal\_evolution.png}}
\end{figure}

\FloatBarrier

\HybridLate\ goes further still: its $T' = 1$ curve already closely tracks the shape of the full
$T' = 50$ curve, oscillating between roughly 63\% and 99\% --- most of the way to final performance
using a single timestep. This is the expected extreme of the same pattern: with only layer 6
spiking, six of the seven layers are instantaneous analog processing, so the vast majority of the
network's rotation-relevant computation is available regardless of the number of spiking timesteps
allowed.

\subsection{A Clean, Reproducible Ordering, and a Compute-Budget Implication}

Combining Sections 6.1--6.3 gives a consistent, monotonic ordering across all three seeds of every
configuration: the more layers that spike, the more timesteps the network needs before its
rotation-robust accuracy structure appears. \FullSpiking\ (7/7 spiking) needs the most accumulation
time, starting from true chance at $T' = 1$; \Alternating\ (3/7) needs less, starting from a partial
signal; \HybridLate\ (1/7) needs almost none, starting from a strong signal already at $T' = 1$. This
ordering is exactly what one would predict from counting the number of spiking ``bottleneck'' stages
a signal must pass through before reaching the classifier, and it held identically in every one of
the nine seed-runs examined (3 configurations $\times$ 3 seeds).

\paragraph{Observation 4 (Compute-budget implication).} \textit{For every configuration, $T' = 5$ is
already close to the $T' = 50$ curve, and $T' = 10$ is nearly indistinguishable from it. Since all
models here were trained and evaluated with the full $T = 50$ throughout, $T = 50$ may be
substantially more compute than necessary. Future training runs --- in particular the still-pending
\HybridEarly\ / \HybridMid\ configurations, and any controlled ablation of the original report's
$T = 30$ setting (Section~\ref{sec:limitations}) --- could plausibly use a much lower $T$ without
materially changing outcomes, at meaningful savings in GPU wall-clock time.}

\section{Finding 4 --- Direct Equivariance Error (Experiment 5.8)}
\label{sec:finding4}

This section covers Experiment 5.8, the direct ground-truth measurement defined in Eq.~\ref{eq:3}
(Section 2.7): $\| f(R_\theta x) - f(x) \| / \| f(x) \|$, computed directly from the network's final
output, with no representation-similarity proxy involved. Unlike Sections~\ref{sec:finding1}--\ref{sec:finding3},
which examine internal layers, this is the most literal possible test of ``does rotating the input
change the answer.''

\subsection{The Expected Sawtooth, Confirmed in Every Configuration}

As explained in Section 2.7, the $p4$-equivariant architecture guarantees exact (0-error) invariance
only at the four group angles $\{0^\circ, 90^\circ, 180^\circ, 270^\circ\}$, with no guarantee
elsewhere. Every plot below confirms this precisely: equivariance error touches (or comes extremely
close to) 0 at exactly the four group angles in every configuration and every seed, and rises to a
local maximum roughly halfway between consecutive group angles (near
$45^\circ, 135^\circ, 225^\circ, 315^\circ$), producing a repeating sawtooth of period $90^\circ$.
This is exactly the behaviour the architecture's mathematical guarantees predict, and its clean
appearance in all twelve seed-runs shown here (4 configurations $\times$ 3 seeds) is a strong sanity
check that the \texttt{P4BatchNorm2d} equivariance fix (Section~\ref{sec:arch}) is functioning
correctly in fully trained models, not merely in isolated unit tests.

\paragraph{Analog}

\begin{figure}[H]
\centering
\includegraphics[width=0.8\textwidth]{analog_seed0/analog_seed0_5.8_equivariance_error.png}
\caption{\Analog, seed 0 --- equivariance error vs.\ angle.
\texttt{analog\_seed0/analog\_seed0\_5.8\_equivariance\_error.png}}
\end{figure}

\begin{figure}[H]
\centering
\includegraphics[width=0.8\textwidth]{analog_seed1/analog_seed1_5.8_equivariance_error.png}
\caption{\Analog, seed 1.
\texttt{analog\_seed1/analog\_seed1\_5.8\_equivariance\_error.png}}
\end{figure}

\begin{figure}[H]
\centering
\includegraphics[width=0.8\textwidth]{analog_seed2/analog_seed2_5.8_equivariance_error.png}
\caption{\Analog, seed 2.
\texttt{analog\_seed2/analog\_seed2\_5.8\_equivariance\_error.png}}
\end{figure}

\FloatBarrier
\paragraph{Alternating}

\begin{figure}[H]
\centering
\includegraphics[width=0.8\textwidth]{alternating_seed0/alternating_seed0_5.8_equivariance_error.png}
\caption{\Alternating, seed 0.
\texttt{alternating\_seed0/alternating\_seed0\_5.8\_equivariance\_error.png}}
\end{figure}

\begin{figure}[H]
\centering
\includegraphics[width=0.8\textwidth]{alternating_seed1/alternating_seed1_5.8_equivariance_error.png}
\caption{\Alternating, seed 1.
\texttt{alternating\_seed1/alternating\_seed1\_5.8\_equivariance\_error.png}}
\end{figure}

\begin{figure}[H]
\centering
\includegraphics[width=0.8\textwidth]{alternating_seed2/alternating_seed2_5.8_equivariance_error.png}
\caption{\Alternating, seed 2.
\texttt{alternating\_seed2/alternating\_seed2\_5.8\_equivariance\_error.png}}
\end{figure}

\FloatBarrier
\paragraph{Full Spiking}

\begin{figure}[H]
\centering
\includegraphics[width=0.8\textwidth]{full_seed0/full_seed0_5.8_equivariance_error.png}
\caption{\FullSpiking, seed 0.
\texttt{full\_seed0/full\_seed0\_5.8\_equivariance\_error.png}}
\end{figure}

\begin{figure}[H]
\centering
\includegraphics[width=0.8\textwidth]{full_seed1/full_seed1_5.8_equivariance_error.png}
\caption{\FullSpiking, seed 1.
\texttt{full\_seed1/full\_seed1\_5.8\_equivariance\_error.png}}
\end{figure}

\begin{figure}[H]
\centering
\includegraphics[width=0.8\textwidth]{full_seed2/full_seed2_5.8_equivariance_error.png}
\caption{\FullSpiking, seed 2.
\texttt{full\_seed2/full\_seed2\_5.8\_equivariance\_error.png}}
\end{figure}

\FloatBarrier
\paragraph{Hybrid Late}

\begin{figure}[H]
\centering
\includegraphics[width=0.8\textwidth]{hybrid_late_seed0/hybrid_late_seed0_5.8_equivariance_error.png}
\caption{\HybridLate, seed 0.
\texttt{hybrid\_late\_seed0/hybrid\_late\_seed0\_5.8\_equivariance\_error.png}}
\end{figure}

\begin{figure}[H]
\centering
\includegraphics[width=0.8\textwidth]{hybrid_late_seed1/hybrid_late_seed1_5.8_equivariance_error.png}
\caption{\HybridLate, seed 1.
\texttt{hybrid\_late\_seed1/hybrid\_late\_seed1\_5.8\_equivariance\_error.png}}
\end{figure}

\begin{figure}[H]
\centering
\includegraphics[width=0.8\textwidth]{hybrid_late_seed2/hybrid_late_seed2_5.8_equivariance_error.png}
\caption{\HybridLate, seed 2.
\texttt{hybrid\_late\_seed2/hybrid\_late\_seed2\_5.8\_equivariance\_error.png}}
\end{figure}

\FloatBarrier

\subsection{Peak Magnitude Does Not Scale Monotonically with Spiking Extent}

If spiking dynamics simply degraded invariance in proportion to how many layers spike, the sawtooth
peak height (worst-case equivariance error, near
$45^\circ, 135^\circ, 225^\circ, 315^\circ$) should increase monotonically from \Analog\ (0 spiking
layers) through \HybridLate\ (1) through \Alternating\ (3) through \FullSpiking\ (7). The data does
not show this.

\Analog\ --- the zero-spiking control --- has the \emph{highest} peak error of all four
configurations, and \FullSpiking\ --- the maximal-spiking configuration --- has the \emph{lowest}.
This is the opposite of the naive ``spiking hurts invariance'' expectation, and is reported here as
an open, reproducible observation rather than an explained result. A plausible but unconfirmed
reading is that the stochastic, rate-coded nature of spiking input, combined with surrogate-gradient
training, may act as an implicit regularizer that reduces the network's worst-case output sensitivity
at the hardest angles, even though (per Finding 1) internal representations are not becoming more
similar by any of the similarity metrics measured. No experiment run so far isolates \emph{why} this
ordering holds.

\begin{table}[H]
\centering
\begin{tabular}{@{}llc@{}}
\toprule
\textbf{Configuration} & \textbf{Spiking layers} & \textbf{Approx. peak equivariance error (mean over seeds)} \\
\midrule
\Analog & 0/7 & $\sim 0.58$--0.59 (highest) \\
\Alternating & 3/7 & $\sim 0.50$--0.60 \\
\HybridLate & 1/7 & $\sim 0.35$--0.40 \\
\FullSpiking & 7/7 & $\sim 0.33$--0.42 (lowest) \\
\bottomrule
\end{tabular}
\caption{Approximate peak equivariance error by configuration, averaged over seeds.}
\label{tab:peakerr}
\end{table}

\subsection{Interpretation Together with Finding 1}

Findings 1 and 4 combine into a genuinely interesting, mildly paradoxical picture. Finding 1 showed
internal representation similarity to the upright case degrades similarly across all four
configurations, bottoming out at layers 6--7 regardless of spiking extent. Finding 4 shows that the
output-level equivariance error not only differs meaningfully between configurations, but differs in
the opposite direction from a naive ``more spiking $\Rightarrow$ worse invariance'' prediction.
Together this means that whatever occurs in the last two layers causing representational drift does
not translate linearly into output error, and spike placement affects the classifier's output-level
behaviour in a way that is not visible from intermediate-layer similarity metrics alone. This is a
concrete argument, for the eventual paper, that both a representation-level analysis (Findings 1--3)
and a direct output-level analysis (Finding 4) are necessary --- research-plan Sections 4 and 5.8
together, neither alone tells the full story.

\section{Synthesis}
\label{sec:synthesis}

The four findings above were run somewhat independently but tell one coherent story when read
together.

\subsection{The Story in Five Statements}

\begin{enumerate}[leftmargin=1.8em]
\item \textbf{Accuracy alone says almost nothing.} Once training is done properly
(\texttt{P4BatchNorm2d} fix, 20 epochs, $T = 50$, 3 seeds; Section~\ref{sec:setup}), all four
configurations reach nearly identical rotation-accuracy curves regardless of how many layers spike,
and the dramatic differences reported earlier were artifacts of the uncorrected setup, not real
properties of spike placement.
\item \textbf{But representation-level analysis (Finding 1) shows the network is not internally
doing the same thing in every configuration.} There is a real, consistent structural pattern ---
invariance holding through layers 1--5, breaking down at 6--7 --- present in all configurations,
including the non-spiking one. This makes it an architectural signature of the $p4$-CNN design, not
a spiking signature.
\item \textbf{Individual-neuron analysis (Finding 2) shows spike placement leaves a fingerprint in
the \emph{shape} of how sensitivity builds up across depth} --- \Alternating\ shows a sharp late-layer
spike, \FullSpiking\ and \Analog\ show a smooth full-depth climb, \HybridLate\ is nearly flat --- even
though the endpoint (layer 7 most sensitive) is common to all. Temporal integration within a layer,
however, does not smooth sensitivity out over time as research-plan Section 9.4 hypothesized;
sensitivity is set within the first few timesteps and then plateaus.
\item \textbf{How quickly rotation-robust accuracy becomes available over spiking time (Finding 3)
scales cleanly and reproducibly with the number of spiking layers} --- more spiking layers means a
longer ``warm-up'' before the W-shaped accuracy curve emerges. This is a clean, mechanistically
explainable result the original report never surfaced, since it never varied $T$ at fixed, correctly
trained models.
\item \textbf{The direct output-level measurement (Finding 4) delivers the most surprising result:}
peak equivariance error is \emph{not} monotonic in the amount of spiking --- \Analog\ has the worst
peak error, \FullSpiking\ the best --- the opposite of a naive ``spiking hurts invariance'' story,
and unexplained pending further work.
\end{enumerate}

\begin{keypoint}
Spike dynamics in this architecture do not determine \emph{whether} rotation-invariant
representations form --- the $p4$-equivariant architecture and its final layers do that, uniformly,
regardless of spiking --- but spike placement does measurably shape \emph{how} that invariance is
approached, both in depth (the per-neuron sensitivity profile across layers) and in time (how many
timesteps are needed before rotation-robust accuracy appears), and, unexpectedly, does not
straightforwardly worsen the network's final output-level equivariance error as spiking is added.
\end{keypoint}

\subsection{Relation to the Research Plan's Theoretical Program}

Research-plan Section 9 proposes a theoretical chain,
\[
\text{Equivariant Spatial Layers} \implies \text{Equivariant Spike Dynamics} \implies
\text{Equivariant Spike Counts} \implies \text{Invariant Classifier},
\]
together with a contraction hypothesis $d_{t+1} \leq \beta d_t$ (Section 9.5 of the research plan)
explaining temporal stabilization. Finding 2 (Section~\ref{sec:54}) provides direct empirical evidence
against the contraction hypothesis as stated: $\mathcal{R}_i(t)$ does not decrease with $t$ in any
configuration tested; it plateaus. This does not invalidate the broader equivariance chain ---
Finding 4's clean sawtooth pattern in every configuration is consistent with the final ``Invariant
Classifier'' link of the chain holding at the group angles --- but it does mean any future
theoretical account (research-plan Sections 9.4--9.5) needs to explain a plateau, not a decay, and
should treat the observed non-monotonic relationship between spiking extent and peak equivariance
error (Finding 4) as a phenomenon requiring its own explanation, separate from the temporal-averaging
story.

\subsection{Relation to Kim et al.\ (2022)}

Kim et al., ``Exploring Temporal Information Dynamics in Spiking Neural Networks'' (NeurIPS 2022),
report a structurally similar decoupling in an unrelated task: their Temporal Information
Concentration (TIC) finding is that Fisher Information about the input concentrates into early
timesteps as training progresses, affecting robustness metrics but not final accuracy. This is the
same qualitative shape as the central finding of this report --- representation/robustness-level
metrics differentiate configurations meaningfully while final accuracy does not --- arrived at
independently, in a different task and metric. This is worth citing directly in the eventual paper's
discussion as external precedent for this kind of decoupling being a real, recurring phenomenon in
SNN research rather than an artifact of this project's particular setup. It also suggests a concrete,
cheap next step (Section~\ref{sec:limitations}): implementing Kim et al.'s Fisher Information trace
$I_t$ and Information Centroid (IC) metric on the existing checkpoints, with no retraining required,
would give a single-number summary directly comparable to their own layer-wise finding (shallow
layers concentrate late-timestep information, deep layers concentrate early-timestep information)
against this report's own layer-6/7 finding.

\section{Limitations and Next Steps}
\label{sec:limitations}

\subsection{Limitations of This Analysis}

\begin{itemize}[leftmargin=1.5em]
\item Only 2 of the 5 spike-placement configurations from the original report's comparison have
multi-seed representation analysis alongside \Analog, \Alternating, and \FullSpiking. \HybridEarly\
and \HybridMid\ have not been trained under the corrected setup, so this report says nothing about
them.
\item No controlled ablation has isolated which of the three candidate factors (\texttt{P4BatchNorm2d}
fix, longer training, higher $T$) is responsible for the original report's dramatic findings not
replicating; all three changed simultaneously between the old and new setups.
\item The counter-intuitive Finding 4 result (\Analog\ worst peak equivariance error, \FullSpiking\
best) has no mechanistic explanation yet; it is reported as an observed, reproducible pattern, not a
settled cause-and-effect claim.
\item SVCCA values should not be read as an absolute 0-to-1 similarity scale (Section~\ref{sec:repsim})
--- the noise floor at the batch sizes used here is $\sim 0.2$--0.4.
\item $\mathcal{R}_i$'s absolute scale is not comparable across spiking vs.\ analog activation types
(Section~\ref{sec:finding2}); only within-configuration comparisons across layers or seeds are on
safe footing.
\item All results use a single fixed task (4-class rotated MNIST, digits 3/4/8/9); generalization to
the full 10-class task or other datasets is untested.
\end{itemize}

\subsection{Recommended Next Steps, in Priority Order}

\begin{enumerate}[leftmargin=1.8em]
\item Train \HybridEarly\ and \HybridMid\ under the corrected setup, and extend Findings 1--4 to
include them, completing the original report's full 5-configuration picture.
\item Run the controlled ablation reproducing the original report's 5-epoch/$T = 30$/CPU regime
inside the corrected codebase, to isolate which factor drove the non-replication of the original
report's two dramatic findings.
\item Implement the Kim et al.\ Fisher Information trace $I_t$ and Information Centroid metric on
existing checkpoints (cheap, no retraining) as a precise complement to Finding 3, and cite Kim et
al.\ directly in the paper's discussion (Section~\ref{sec:synthesis}).
\item Given Finding 3's compute-budget observation ($T' = 10$ already nearly matches $T' = 50$ for
every configuration tested), consider reducing $T$ for the remaining \HybridEarly/\HybridMid\
training runs and the ablation study above, to save GPU time on \texttt{gpuserver}.
\item Cheap, scoped additions when time allows: latent trajectory PCA plots ($0^\circ$ vs.\
$90^\circ$ paths through $h_\ell(t, x)$) using the existing temporal-representation extraction
infrastructure with no retraining; a narrowly scoped check of whether surrogate-gradient training
keeps the equivariant weight tensor's gradients consistent under rotation (research-plan Section
9.1).
\item Update the standalone internship report to incorporate this multi-seed representation analysis
in place of the earlier seed-0-only version, and fold in the Kim et al.\ discussion once Step 3 is
complete.
\end{enumerate}

\subsection{Explicitly Deprioritized (Not Recommended Before the Deadline)}

\begin{itemize}[leftmargin=1.5em]
\item Full PLIF (learnable membrane decay, research-plan Section 5.6) retraining across all
configurations --- real and cheap to add to the codebase, but requires retraining every
configuration again.
\item DVS/N-MNIST spatio-temporal group equivariance extension --- a substantially larger, separate
project requiring a new dataset, new encoding, and new group-theoretic construction; not realistic
before the ICLR abstract/full-paper deadlines.
\item Full gated-recurrent-cell (GRU/LSTM-style) architecture extension combined with
$p4$-equivariance --- genuinely novel but a multi-week undertaking requiring the same
equivariance-verification rigor as the \texttt{P4BatchNorm2d} fix.
\item Rotation-aware pruning (research-plan Section 5.5), learnable-threshold experiments beyond
PLIF, a reduced-parameter-count architecture variant, and the 10-class task extension --- all
previously proposed, none run, all lower priority than completing the core 5-configuration
comparison.
\end{itemize}

\end{document}
